% GENERATED FILE. Edit canonical lessons, then run npm run generate:books.
% canonical-source-hash: fnv1a64:57cd4f6f50f3afbe
% canonical-lessons: LA-C336-gladiator, LA-C336-lanista, LA-C336-munus, LA-C336-venatio, LA-C336-bestia

\chapter{Mūnus}
\label{ch:la-munus}
\hlchaptermodality{336}

\begin{chapteropening}
\textbf{By the end of this chapter:} \emph{I can describe the shows in the arena: the gladiator, the trainer, the public show, the wild-beast hunt and the beast.}

Complete the last lesson of chapter 336: I can describe the shows in the arena: the gladiator, the trainer, the public show, the wild-beast hunt and the beast.
\end{chapteropening}

\section[gladiātor, gladiātōris]{gladiātor, gladiātōris — a gladiator}

\label{lesson:LA-C336-gladiator}



\begin{quote}
Before the new one: say the Latin for a racing team, then the Latin for of the circus.
\end{quote}

\subsection*{You'll want to know: gladiātor, gladiātōris}
\textbf{gladiātor, gladiātōris} — \textquotedblleft{}a gladiator\textquotedblright{}.

\begin{quote}
Gladiātōrēs in harēnā pugnābant.
\end{quote}

\begin{quote}
\emph{Gladiators were fighting in the arena.}
\end{quote}

Masculine, 3rd declension, from \emph{gladius}, 'sword' (taught). Most were slaves or prisoners; a few were volunteers.

\begin{cousinweb}
\textbf{gladiātor} is built on \emph{gladius}, \textquotedblleft{}sword\textquotedblright{} (EtymDB): the swordsman. English \emph{gladiator} is the same word.
\end{cousinweb}

\begin{grammarlens}[title={What you've built}]
One more word for telling, explaining and describing at B1.
\end{grammarlens}

\subsection*{Your turn}
\begin{itemize}
\raggedright
  \item \emph{Say it:} \emph{gladiātor, gladiātōris}
  \item \emph{Say it:} \emph{gladiātor, gladiātōris}, once more
  \item \emph{Say it:} a sentence of your own with \emph{gladiātor, gladiātōris}
  \item \emph{Recall:} say the Latin for a racing team, then the Latin for of the circus, then say \emph{gladiātor, gladiātōris} again
\end{itemize}

\begin{tcolorbox}[breakable,colback=teal!4,colframe=teal!35!black,title={Before you move on}]
What does gladiātor, gladiātōris mean? (\textquotedblleft{}A gladiator\textquotedblright{}.) Say it once more.
\end{tcolorbox}
\section[lanista, lanistae]{lanista, lanistae — a trainer of gladiators}

\label{lesson:LA-C336-lanista}



\begin{quote}
Before the new one: say the Latin for of the circus, then the Latin for a gladiator.
\end{quote}

\subsection*{You'll want to know: lanista, lanistae}
\textbf{lanista, lanistae} — \textquotedblleft{}a trainer of gladiators\textquotedblright{}.

\begin{quote}
Lanista novōs gladiātōrēs ēmit.
\end{quote}

\begin{quote}
\emph{The trainer bought new gladiators.}
\end{quote}

Masculine, 1st declension. He owned and trained a troop; in Cicero the word is an insult.

\begin{cousinweb}
\textbf{lanista} is a word whose origin is not settled (Lewis \& Short).
\end{cousinweb}

\begin{grammarlens}[title={What you've built}]
One more word for telling, explaining and describing at B1.
\end{grammarlens}

\subsection*{Your turn}
\begin{itemize}
\raggedright
  \item \emph{Say it:} \emph{lanista, lanistae}
  \item \emph{Say it:} \emph{lanista, lanistae}, once more
  \item \emph{Say it:} a sentence of your own with \emph{lanista, lanistae}
  \item \emph{Recall:} say the Latin for of the circus, then the Latin for a gladiator, then say \emph{lanista, lanistae} again
\end{itemize}

\begin{tcolorbox}[breakable,colback=teal!4,colframe=teal!35!black,title={Before you move on}]
What does lanista, lanistae mean? (\textquotedblleft{}A trainer of gladiators\textquotedblright{}.) Say it once more.
\end{tcolorbox}
\section[mūnus, mūneris]{mūnus, mūneris — a public show; a gift; a duty}

\label{lesson:LA-C336-munus}



\begin{quote}
Before the new one: say the Latin for a gladiator, then the Latin for a trainer of gladiators.
\end{quote}

\subsection*{You'll want to know: mūnus, mūneris}
\textbf{mūnus, mūneris} — \textquotedblleft{}a public show; a gift; a duty\textquotedblright{}.

\begin{quote}
Ad mūnus tōta urbs vēnit.
\end{quote}

\begin{quote}
\emph{The whole city came to the show.}
\end{quote}

Neuter, 3rd declension. First a duty, then a gift, and so the gladiator show given to the people. English \emph{munificent}.

\begin{cousinweb}
\textbf{mūnus} is inherited from Proto-Italic \emph{*moinos}, \textquotedblleft{}service\textquotedblright{} (EtymDB): a duty, then a gift, then a show. English \emph{munificent} comes from its family.
\end{cousinweb}

\begin{grammarlens}[title={What you've built}]
One more word for telling, explaining and describing at B1.
\end{grammarlens}

\subsection*{Your turn}
\begin{itemize}
\raggedright
  \item \emph{Say it:} \emph{mūnus, mūneris}
  \item \emph{Say it:} \emph{mūnus, mūneris}, once more
  \item \emph{Say it:} a sentence of your own with \emph{mūnus, mūneris}
  \item \emph{Recall:} say the Latin for a gladiator, then the Latin for a trainer of gladiators, then say \emph{mūnus, mūneris} again
\end{itemize}

\begin{tcolorbox}[breakable,colback=teal!4,colframe=teal!35!black,title={Before you move on}]
What does mūnus, mūneris mean? (\textquotedblleft{}A public show; a gift; a duty\textquotedblright{}.) Say it once more.
\end{tcolorbox}
\section[vēnātiō, vēnātiōnis]{vēnātiō, vēnātiōnis — a hunt; a wild-beast show}

\label{lesson:LA-C336-venatio}



\begin{quote}
Before the new one: say the Latin for a trainer of gladiators, then the Latin for a public show.
\end{quote}

\subsection*{You'll want to know: vēnātiō, vēnātiōnis}
\textbf{vēnātiō, vēnātiōnis} — \textquotedblleft{}a hunt; a wild-beast show\textquotedblright{}.

\begin{quote}
Māne vēnātiō erat, post merīdiem gladiātōrēs.
\end{quote}

\begin{quote}
\emph{In the morning there was a wild-beast hunt, in the afternoon the gladiators.}
\end{quote}

Feminine, from \emph{vēnor}, 'to hunt' (taught). In the arena, hunters fought lions and bears.

\begin{cousinweb}
\textbf{vēnātiō} is built on \emph{vēnor}, \textquotedblleft{}to hunt\textquotedblright{} (Lewis \& Short): hunting, and the meat of the hunt. English \emph{venison} comes from it through French.
\end{cousinweb}

\begin{grammarlens}[title={What you've built}]
One more word for telling, explaining and describing at B1.
\end{grammarlens}

\subsection*{Your turn}
\begin{itemize}
\raggedright
  \item \emph{Say it:} \emph{vēnātiō, vēnātiōnis}
  \item \emph{Say it:} \emph{vēnātiō, vēnātiōnis}, once more
  \item \emph{Say it:} a sentence of your own with \emph{vēnātiō, vēnātiōnis}
  \item \emph{Recall:} say the Latin for a trainer of gladiators, then the Latin for a public show, then say \emph{vēnātiō, vēnātiōnis} again
\end{itemize}

\begin{tcolorbox}[breakable,colback=teal!4,colframe=teal!35!black,title={Before you move on}]
What does vēnātiō, vēnātiōnis mean? (\textquotedblleft{}A hunt; a wild-beast show\textquotedblright{}.) Say it once more.
\end{tcolorbox}
\section[bēstia, bēstiae]{bēstia, bēstiae — a beast, a wild animal}

\label{lesson:LA-C336-bestia}



\begin{quote}
Before the new one: say the Latin for a public show, then the Latin for a hunt.
\end{quote}

\subsection*{You'll want to know: bēstia, bēstiae}
\textbf{bēstia, bēstiae} — \textquotedblleft{}a beast, a wild animal\textquotedblright{}.

\begin{quote}
Bēstiae ingentēs in harēnam missae sunt.
\end{quote}

\begin{quote}
\emph{Huge beasts were sent into the arena.}
\end{quote}

Feminine. A creature without reason, as opposed to people. English \emph{beast}, through French.

\begin{cousinweb}
\textbf{bēstia} is a Latin word whose origin is not settled; Lewis \& Short suggest a link with \emph{fera}, \textquotedblleft{}wild beast\textquotedblright{}. English \emph{beast} comes from it through French.
\end{cousinweb}

\begin{grammarlens}[title={What you've built}]
One more word for telling, explaining and describing at B1.
\end{grammarlens}

\subsection*{Your turn}
\begin{itemize}
\raggedright
  \item \emph{Say it:} \emph{bēstia, bēstiae}
  \item \emph{Say it:} \emph{bēstia, bēstiae}, once more
  \item \emph{Say it:} a sentence of your own with \emph{bēstia, bēstiae}
  \item \emph{Recall:} say the Latin for a public show, then the Latin for a hunt, then say \emph{bēstia, bēstiae} again
\end{itemize}

\begin{tcolorbox}[breakable,colback=teal!4,colframe=teal!35!black,title={Before you move on}]
What does bēstia, bēstiae mean? (\textquotedblleft{}A beast, a wild animal\textquotedblright{}.) Say it once more.
\end{tcolorbox}
